\newpage
\textbf{\huge CHAPTER  2}\\
\section{LITERATURE REVIEW} \setcounter{section}{2}
\subsection{Overview of Distributed Stream Processing}\setcounter{subsection}{1}

\begin{justify}
    The field of distributed stream processing has evolved dramatically over the past two decades, driven by the explosion of internet-scale data and the increasing demand for low-latency analytics. Traditional batch processing paradigms, such as Hadoop MapReduce \cite{[1]}, were designed to process large datasets at rest, making them fundamentally unsuitable for applications that require near real-time responsiveness. This limitation catalysed the development of dedicated stream processing frameworks capable of handling continuous, unbounded data streams with high throughput and fault tolerance.

    Early stream processing systems such as Aurora \cite{[2]} and STREAM \cite{[3]} introduced the theoretical foundations for continuous query processing over data streams. However, these systems lacked the distributed scalability required to handle the data volumes generated by modern web-scale platforms. The advent of frameworks like Apache Storm, Apache Samza, Apache Spark Streaming, and ultimately Apache Flink marked a new era of distributed, fault-tolerant, and scalable stream processing \cite{[4]}.
\end{justify}

\vspace{0.25cm}

\subsection{Apache Kafka: The Distributed Messaging Backbone}
\begin{justify}
    Apache Kafka, originally developed at LinkedIn and open-sourced in 2011, has become the de facto standard for high-throughput, fault-tolerant distributed messaging in modern data architectures \cite{[5]}. Kafka is built around a distributed commit log abstraction: producers publish records to named \textit{topics}, and consumers pull records from these topics at their own pace. Topics are partitioned and replicated across a cluster of broker nodes, providing both horizontal scalability and fault tolerance.

    Kafka's design philosophy centres on durability and decoupling. Its append-only log model guarantees that records are persisted for a configurable retention period, allowing consumers to replay events from any historical offset. This characteristic is particularly valuable in stream processing architectures, where exactly-once delivery requires the ability to rewind and reprocess events after a failure recovery. Kafka's consumer group model, where each partition is consumed by exactly one consumer within a group, enables parallel, ordered processing across multiple stream processor instances \cite{[6]}.

    In enterprise-scale deployments like Flipkart, Kafka serves as the central nervous system, decoupling upstream business services (catalog management, pricing engines, inventory systems) from downstream consumers (search ingestion pipelines, analytics systems, ML platforms). Its ability to sustain millions of messages per second with sub-10ms publish latency makes it uniquely suited for real-time e-commerce use cases.
\end{justify}

\vspace{0.25cm}

\subsection{Apache Storm: Legacy Real-Time Processing}
\begin{justify}
    Apache Storm, created by Nathan Marz at BackType and later acquired by Twitter before being donated to the Apache Software Foundation in 2011, was one of the first widely adopted open-source distributed stream processing systems \cite{[7]}. Storm introduced the concept of a \textit{Topology} — a directed acyclic graph (DAG) of computation composed of \textit{Spouts} (data sources) and \textit{Bolts} (processing units). Tuples flow through this topology, with each bolt performing discrete transformations such as filtering, aggregation, or external lookups.

    Storm's design prioritises simplicity and low-latency tuple processing, offering \textit{at-least-once} processing guarantees through its ACK-based tracking mechanism. Each tuple is tracked from its origin spout through the entire bolt chain; if any bolt fails to acknowledge a tuple within a timeout window, Storm automatically re-emits the tuple from the spout. While this ensures no data loss, it does not prevent duplicate processing — downstream systems must be designed to be idempotent.

    A fundamental limitation of Apache Storm is its lack of native state management. Storm is essentially a stateless processing engine: bolts compute results tuple-by-tuple without any built-in mechanism to maintain and access processing state across events. Consequently, any stateful operation—such as maintaining the current composite view of a product entity—requires bolts to issue synchronous read and write calls to an external state store (typically Apache HBase or Redis). This external I/O dependency introduces significant network latency, limits throughput, and creates tight coupling between the processing layer and the storage layer. Studies have documented that at high throughput, Storm's external state access can become the dominant source of end-to-end latency \cite{[8]}.

    Additionally, Storm's event-time handling capabilities are rudimentary. Storm processes events in \textit{processing time} by default, meaning the order of processing depends on the order of arrival rather than the logical event timestamp. This makes it challenging to correctly handle out-of-order events that are common in distributed, asynchronous environments — a significant drawback for applications requiring strict temporal consistency.
\end{justify}

\vspace{0.25cm}

\subsection{Apache Flink: Modern Stateful Stream Processing}
\begin{justify}
    Apache Flink, originally developed as the Stratosphere research project at TU Berlin and donated to the Apache Software Foundation in 2014, represents the current state-of-the-art in distributed stream processing \cite{[9]}. Flink's foundational design principle — ``state as a first-class citizen'' — directly addresses the core limitations of its predecessors.

    \textbf{Unified DataStream API.} Flink provides a rich, expressive DataStream API for building complex stream processing pipelines. Operations such as \textit{map}, \textit{filter}, \textit{flatMap}, \textit{keyBy}, \textit{reduce}, and \textit{process} can be composed into sophisticated transformation graphs. The \textit{keyBy()} operator is particularly powerful: it partitions the stream by a user-defined key, guaranteeing that all events with the same key are routed to the same parallel subtask, enabling correct, co-located stateful computations without inter-node coordination \cite{[10]}.

    \textbf{Managed State and RocksDB Backend.} Flink natively manages operator state and keyed state within the processing runtime. The \textit{RocksDB} state backend stores state on local disk (using a log-structured merge-tree), enabling state sizes that far exceed available memory while maintaining sub-millisecond access times. This eliminates the need for external state stores during active stream processing — a fundamental departure from Storm's architecture \cite{[11]}.

    \textbf{Event Time and Watermarks.} Flink provides first-class support for \textit{event time} processing, where the progress of computation is governed by timestamps embedded in the data rather than wall-clock time. \textit{Watermarks} serve as progress markers that inform the runtime of the maximum event timestamp seen so far (minus an allowed out-of-order tolerance), enabling correct handling of late-arriving events without compromising throughput \cite{[12]}.

    \textbf{Exactly-Once Semantics via Chandy-Lamport Checkpointing.} Flink's fault tolerance mechanism is based on the Chandy-Lamport distributed snapshot algorithm. Periodic \textit{checkpoint barriers} are injected into the data stream; as these barriers propagate through the operator graph, each operator asynchronously snapshots its state to a durable storage backend (e.g., HDFS). Upon failure, Flink restores the entire operator graph from the latest consistent checkpoint and rewinds source offsets (e.g., Kafka consumer offsets), guaranteeing exactly-once processing semantics end-to-end \cite{[9]}.
\end{justify}

\vspace{0.25cm}

\subsection{Comparison of Stream Processing Frameworks}
\begin{justify}
    Table~\ref{tab:framework_comparison} presents a structured comparison of the three major stream processing frameworks relevant to this project: Apache Storm, Apache Spark Streaming (a widely used alternative), and Apache Flink.

    \begin{table}[h]
        \centering
        \small
        \caption{Comparison of Stream Processing Frameworks}
        \renewcommand{\arraystretch}{1.5}
        \begin{tabular}{|p{3.5cm}|p{3cm}|p{3.5cm}|p{3.5cm}|}
            \hline
            \textbf{Feature} & \textbf{Apache Storm} & \textbf{Spark Streaming} & \textbf{Apache Flink} \\
            \hline
            Processing Model & Record-at-a-time & Micro-batch & Record-at-a-time \\
            \hline
            State Management & External only & RDD checkpointing & Native (RocksDB) \\
            \hline
            Processing Semantics & At-least-once & Exactly-once & Exactly-once \\
            \hline
            Event Time Support & Limited & Limited & Full (Watermarks) \\
            \hline
            Latency & Sub-second & Seconds (batch interval) & Sub-second \\
            \hline
            Backpressure & Manual & Auto (blocking) & Auto (credit-based) \\
            \hline
            Kafka Integration & Kafka Spout & Kafka DStream & FlinkKafkaConsumer \\
            \hline
            Fault Tolerance & ACK-based (at-least-once) & RDD lineage & Chandy-Lamport \\
            \hline
        \end{tabular}
        \label{tab:framework_comparison}
    \end{table}

    Spark Streaming, while providing exactly-once semantics, operates on a micro-batch model: it buffers incoming events for a configurable interval (typically 500ms--1s) before processing them as a mini-batch. This introduces inherent latency that makes it unsuitable for applications requiring sub-second ingestion, such as Flipkart's search index updates. Flink's record-at-a-time model, combined with its native state management and exactly-once guarantees, makes it the superior choice for latency-sensitive, stateful stream processing at scale \cite{[13]}.
\end{justify}

\vspace{0.25cm}

\subsection{Real-Time Search Data Ingestion: Related Work}
\begin{justify}
    The challenge of maintaining real-time search indices over continuously evolving product catalogs has been extensively studied in the context of large-scale e-commerce and web search systems.

    LinkedIn's Espresso \cite{[14]} pioneered the use of a change-data-capture (CDC) approach, where mutations to a primary database are streamed as events into an indexing pipeline. This pattern is architecturally similar to Flipkart's ingestion system, where upstream business services emit delta signals into Kafka topics. LinkedIn subsequently migrated their stream processing from a custom system to Apache Samza and later to Flink, reporting significant improvements in latency and operational complexity \cite{[5]}.

    Uber's stream processing infrastructure \cite{[15]} demonstrates the industry-wide trend towards Flink adoption for exactly-once, low-latency data pipelines. Uber's migration from Storm to Flink for their real-time analytics platform showed a reduction in processing latency from the order of seconds to milliseconds, validating the architectural advantages of Flink's native state management over Storm's external-state model.

    Amazon's product search indexing system \cite{[1]} employs a pipeline architecture conceptually similar to Flipkart's: signals from catalog, pricing, and inventory systems are ingested via a distributed message bus and processed by a stream engine before being reflected in the live search index. The emphasis on sub-second \textit{time-to-index} (TTI) — the latency between a catalog change event and its visibility in search results — is a universal KPI across major e-commerce platforms, underscoring the importance of low-latency stream processing architectures.

    Collectively, this body of work validates the strategic direction undertaken in this internship: migrating from a stateless, at-least-once Storm topology to a stateful, exactly-once Flink pipeline is a well-established, industry-proven approach to achieving sub-second search index freshness at e-commerce scale.
\end{justify}
